\section{Learning uniform conflict transversals}
\label{sec:euler-transversal}

The smallest conflict in one returned relaxation point need not identify a
small set of tetrahedra that resolves the whole search.  This section gives
an exact uniform parameter for the fixed-source Q cone and an optional
search policy that learns constraints on that parameter.  The policy can
reuse the mask oracle and negative certificate format from
Section~\ref{sec:euler-branch-tree}.  Its general quasi-polynomial bound
remains conditional on a structural theorem and a polynomial LP backend.

\subsection{Coordinate faces retain the root extreme rays}

Fix a finite source triangulation with $t$ tetrahedra and its full
quadrilateral matching cone
\[
 K=\{q\in\mathbb R^{3t}:q\ge0,\ Cq=0\}.
\]
The root may already have a partial mask; in that case replace $K$ by
that coordinate face throughout.  For an allowed mask $J$, zero-extension
identifies the node cone with
\[
 K_J=K\cap\{q_j=0:j\notin J\}.
\]
The fixed-source Euler function is the canonical formal extension
\[
 F(q)=\max_{\alpha\in\mathcal A}c_\alpha^{\mathsf T}q.
\]
It is convex and positively homogeneous.  The number of source anchor
tuples is $A=|\mathcal A|$; after one-boundary-vertex preprocessing,
$A\le4t$.  Incompatible coordinates remain formal matching vectors, not
embedded normal surfaces.

\begin{lemma}[Inheritance under coordinate restriction]
\label{lem:transversal-face-rays}
Every extreme ray of $K_J$ is an extreme ray of $K$.  Moreover, if
$F(q)>0$ anywhere in $K_J$, some extreme ray of $K_J$ has positive $F$.
\end{lemma}
\begin{proof}
Suppose a ray vector $r\in K_J$ decomposes as $u+v$ with $u,v\in K$.
For $j\notin J$, nonnegativity and $u_j+v_j=r_j=0$ force
$u_j=v_j=0$.  Thus $u,v\in K_J$, so extremality in $K_J$ forces both
summands onto the ray of $r$.  This proves the first assertion.

The nonnegative cone is pointed.  Its section $\sum_jq_j=1$ is a
compact polytope when it is nonempty, and its vertices correspond to
extreme rays.  Normalize a positive point into this section and express
it as a convex combination of vertices.  Convexity of $F$ implies that
one of those vertices has positive $F$.
\end{proof}

The source model is important here.  A node restricts columns of the same
cone; rebuilding an equivalent eliminated representation does not create
new geometric extreme rays.  The exact LP oracle returns a nonzero basic
feasible point of $C_Jq=0$, $q\ge0$, $\sum q\le1$.  Its slack variables
are uniquely determined by $q$, so that point projects to an extreme ray
of $K_J$, and hence of the root.  The positive canonical standard lift
may still require the existing full-standard support extraction before
crushing.

\subsection{A uniform parameter and a distinction for nonextreme points}

For a nonnegative Q vector define its complete conflict set by
\[
 D(q)=\{\tau:\#\{i:q_{\tau,i}>0\}\ge2\}.
\]
Let $\mathcal R_+$ be the root extreme rays with $F>0$ that are
incompatible.  Define the hypergraph and its transversal number by
\[
 \mathcal H(T)=\{D(r):r\in\mathcal R_+\},\qquad
 h(T)=\min\{|B|:B\cap D\ne\varnothing
                       \text{ for every }D\in\mathcal H(T)\}.
\]
The empty hypergraph has transversal number zero.  Compatible positive
rays are successful witnesses and are excluded from this hypergraph.
Every remaining edge is nonempty, so $h(T)\le t$.

A $B$-assignment retains one Q type in each tetrahedron of $B$ and all
previously allowed types elsewhere.  Its zero-Q possibility is included
in every singleton choice.

\begin{proposition}[Uniform conflict-transversal characterization]
\label{prop:uniform-q-transversal}
A set $B$ meets every edge of $\mathcal H(T)$ if and only if no
$B$-assignment face contains an incompatible positive root extreme ray.
If there is no admissible canonical positive vector, these conditions
are also equivalent to $F\le0$ on every $B$-assignment face.
\end{proposition}
\begin{proof}
If $B$ meets $D(r)$, resolving a tetrahedron in their intersection
excludes $r$ under every singleton choice.  Conversely, if $B$ misses
$D(r)$, at each tetrahedron of $B$ choose the sole type used by $r$,
choosing arbitrarily when its Q coordinates vanish.  That assignment
retains $r$.  On a negative source every positive extreme ray is
incompatible.  Lemma~\ref{lem:transversal-face-rays} then supplies the
last equivalence.
\end{proof}

\begin{lemma}[A nonextreme positive point still obstructs negative coverage]
\label{lem:nonextreme-conflict-constraint}
Let $q\in K$ have $F(q)>0$, without any extremality hypothesis.  Every
set $B$ for which \emph{all} $B$-assignment relaxations are nonpositive
must meet $D(q)$.
\end{lemma}
\begin{proof}
If $B$ misses $D(q)$, choose the used type at each tetrahedron of $B$,
arbitrarily where all types vanish.  The resulting assignment retains
$q$, contradicting its claimed nonpositivity.
\end{proof}

This lemma gives a valid hitting constraint for complete negative
coverage even from a nonextreme point.  On a negative source,
Proposition~\ref{prop:uniform-q-transversal} identifies such covering
sets with the root-ray transversals.  On a positive source the distinction
matters: an incompatible sum of compatible positive rays need not give
a necessary constraint on a transversal of \emph{incompatible extreme}
rays.  The parameter theorem for both outcomes therefore uses the actual
oracle's basic-ray guarantee.  The finite lower-bound witnesses described
by the research interface are separately checked for Q-support nullity
one; their extremality is not inferred from their formal Euler value.

A bound on one observed $|D(q)|$ is not a bound on $h(T)$.  For example,
$m$ verified positive root rays with pairwise disjoint conflict sets force
$h(T)\ge m$, regardless of the size of any individual set.  Conversely,
a complete negative tree that splits only in $B$ proves $h(T)\le|B|$:
a positive ray avoiding conflicts in $B$ could follow its retained
choices through that tree and would contradict a negative leaf.

\subsection{The optional learning policy}

The maintained \path{search_euler_transversal} interface uses the shared
\path{_EulerMaskOracle}, with the fixed source model and the same bounded
cache of at most $256$ source-row dual bounds.  It records the complete
conflict set of every incompatible positive basic ray.  It then explores
a candidate $B$ using ordinary complete mask branching, but permits
splits only at returned conflict tetrahedra in $B$.

If a returned conflict set is disjoint from $B$, it refutes that
candidate.  The partial candidate tree is discarded; all source algebra,
checked dual bounds, learned conflicts and charged work are retained.
A compatible positive vector is checked against original normal matching
and returned immediately.  A successful negative candidate emits the
existing \path{normal-euler-branch-tree-v1} packet with source-model v2
negative leaves.  Failed candidates never become negative evidence.

Candidates are generated by one stream of subsets, first by cardinality
and then lexicographically.  Every subset is tested against the currently
known hypergraph.  A skipped subset remains invalid because constraints
only accumulate.  An explored candidate that fails acquires a disjoint
conflict and also remains invalid.  Therefore continuing the stream
still selects a minimum hitting set of the known hypergraph, without
retesting earlier subsets.

\begin{proposition}[Conditional learning bound]
\label{thm:transversal-learning-bound}
Disable operational caps and assume an exact polynomial-bit LP procedure
with polynomial basic-ray recovery.  The learning policy is complete.
Writing $h=h(T)$, it tests at most
\[
 \sum_{k=0}^{h}\binom tk
\]
subsets, and its total oracle-node count is at most
\[
 \sum_{k=0}^{h}\binom tk\frac{3^{k+1}-1}{2}.
\]
Its bit work is $A(3t)^{O(h)}\operatorname{poly}(t)$.  A supplied uniform transversal
of size $h$ needs only $A3^{O(h)}\operatorname{poly}(t)$ work.
\end{proposition}
\begin{proof}
Every learned edge is the conflict set of an incompatible positive
root extreme ray, by Lemma~\ref{lem:transversal-face-rays}.  A true
size-$h$ transversal therefore meets all learned edges.  No attempted
minimum candidate can have more than $h$ elements.  The monotone stream
tests each subset at most once.

A candidate of size $k$ has a tree of depth at most $k$, with at most
three children at each split.  Any incompatible ray that survives a
resolved candidate gives a new disjoint conflict and rejects it.  For
a true transversal this cannot happen: every positive basic ray in
an assignment face is compatible.  Hence that candidate returns a
positive witness or a complete negative tree.  The finite source has
only finitely many subsets, and the full ambiguous set always resolves
all conflicts, proving completeness independently of a small parameter.

Each oracle node uses at most one envelope query and $A$ anchor queries,
with polynomial source arithmetic and proof data under the stated LP
hypothesis.  Hypergraph membership tests, subset generation and retained
conflicts add a polynomial in the number of charged nodes and subsets,
which remains $(3t)^{O(h)}\operatorname{poly}(t)$.  The supplied-set statement removes
the subset-learning factor.
\end{proof}

The native producer retains exact Bland simplex and makes no polynomial
pivot claim.  Its node, pivot and anchor allowances are shared across
all candidate restarts.  A separate \path{max_hitting_subsets} allowance,
$10{,}000$ by default, charges every tested subset, including those rejected
before an LP call.  All caller interruptions propagate, and any exhausted
allowance returns \path{INCONCLUSIVE} without a partial negative packet.
The diagram driver exposes the policy explicitly through
\path{sector_method='transversal'} and keeps the source-bound crushing
continuation and independent global replay.

Without a small $h$ theorem, the crude sum over all candidate trees is
$O(4^t)$ oracle nodes, compared with $O(3^t)$ for one ordinary full tree.
This policy is consequently optional.  A possible later default could
alternate fresh deterministic DFS and learning attempts with geometrically
increasing node allowances.  Summing those allowances gives constant-factor
competitiveness in charged oracle nodes, with polynomial preparation
repeated logarithmically many times.  This is not automatically a bound
on elapsed time: pivot, anchor and subset-selection work must also be
scheduled or capped so that an expensive inner operation cannot starve
the other policy.  Such a dovetail is not part of this implementation.

\subsection{Why local uniform-Q normalization does not close the gap}

There is a useful local simplification, but positive Q basic rays already
satisfy it.  Assume now the finite triangulation has exactly one boundary
vertex with disc link $L$, so $\chi(L)=1$.  For each tetrahedron $\tau$,
let $W_\tau$ have $+1$ in its four triangle coordinates, $-1$ in its
three Q coordinates, and zero elsewhere.  Each face-arc expression is
triangle plus quad, so $W_\tau$ satisfies the original matching equations.
Its face-arc and edge-weight contributions vanish and its disc contribution
is $4-3=1$; hence $\chi(W_\tau)=1$.

These are the tetrahedral matching solutions in the Kang--Rubinstein
basis theorem, in the finite-boundary formulation of Luo--Tillmann.\footnote{%
Feng Luo and Stephan Tillmann, \emph{Angle structures and normal surfaces},
\url{https://arxiv.org/pdf/math/0510537v2}, Theorem~6 and Lemma~7
on printed page~5, and the Euler calculation on printed page~6.
The basis theorem explicitly allows compact pseudomanifolds with boundary.}
Let $r_\tau$ be the Q vector $(1,1,1)$ in $\tau$ and zero elsewhere.
It is the Q projection of $-W_\tau$, so $Cr_\tau=0$.  Its canonical
nonnegative lift is $L-W_\tau$, including when $t=1$.  Therefore
\[
 F(r_\tau)=0.
\]
For arbitrary $q\in K$, subtract
$m_\tau=\min_iq_{\tau,i}$ copies of each $r_\tau$.  The resulting $q'$
remains in $K$, has a zero Q coordinate in every tetrahedron, and
\[
 F(q)\le F(q')+\sum_\tau m_\tau F(r_\tau)=F(q')
\]
by subadditivity.  Admissible vectors are unchanged.  Equivalently, if
$p$ are the canonical triangle coordinates of $q$ and
$\mu=\min_{\tau,v}(p_{\tau,v}+m_\tau)$, its exact Euler gain is
$\sum_\tau m_\tau-\mu\ge0$.

If an extreme ray has all three types positive in $\tau$, subtraction
of its positive minimum decomposes it into $m_\tau r_\tau$ and another
nonnegative matching vector.  Extremality forces it onto $r_\tau$, where
$F=0$.  Thus every genuinely positive Q extreme ray already has a zero
Q entry in every tetrahedron.  Enforcing that condition cannot improve
the positive-existence oracle.  Its three two-type faces are overlapping
relaxations and do not replace the three admissible singleton choices;
the remaining two-type conflicts still require resolution.

Nor does removing these directions give a polylogarithmic linear space.
The one-vertex boundary torus has two triangles.  Face incidence and
$\chi(M)=0$ give $e=t+2$ global edges.  The cited basis theorem gives
standard matching dimension $t+e=2t+2$.  The triangle-only kernel has
dimension one, so Q matching has dimension $2t+1$.  Quotienting its
linear space by the $t$ independent vectors $r_\tau$ leaves dimension
$t+1$.  This is a linear-algebra statement, not a claim that $F$ descends
to that quotient or a lower bound on every search tree.

The next global obligation is therefore a bound on a uniform resolving
parameter, or a different structured search theorem, rather than the
already automatic minimum-Q normalization.  A polylogarithmic bound on
$h$ at all authenticated one-vertex stages, together with the polynomial
LP hypothesis, would give quasi-polynomial search through the polynomially
many crushing stages.  No such universal bound is established here.

\subsection{Fresh finite audit and exact resolving numbers}

The source-pinned correctness run in
\texttt{data/euler-transversal-audit.json} exercises the maintained native
policy, with a separate replay in
\texttt{data/euler-transversal-replay.txt}.  The replay command is
\texttt{python -B -S -m euler\_transversal\_research.check\_evidence};
it imports neither Regina nor an external optimization package and performs
no new surface discovery or timing.  Pins cover the entire native Python
package and the audit's research dependencies.  This experiment establishes
finite agreement and exact certificates, not an asymptotic bound.

The supplied-source audit freshly enumerates the original \emph{standard}
normal vertex surfaces using Regina: 35 vertices over three sources.  It
compares all $8^2=64$ Q masks of the two-tetrahedron layered torus, and twelve
masks each of sources with an interior vertex and multiple boundary vertices.
Both envelope-first and anchor-only learning run on all 88 masks.  The 31
positive masks give 62 positive queries; the 57 negative masks give 114
negative queries.  Every positive result has original matching and canonical
Euler checked independently, then undergoes native full-standard-ray
extraction.  Its extracted coordinates occur in the freshly enumerated
original standard list, and its labelled crushing quotient agrees exactly
with Regina in all 62 comparisons.  Every negative query carries a complete
independently replayed branch tree.  The saved-list replay verifies each
listed ray and each comparison; completeness of those lists is the evidence
from fresh Regina enumeration, not a new enumeration by the replay checker.

The same audit constructs three exteriors from their diagrams, including
shelling, vertex reduction, all search work, and complete source-bound
recognition proofs.  The crossing-free circle is certified \texttt{UNKNOT};
the braid $(2;1,1,1)$ and the braid $(3;1,-2,1,-2)$ are certified
\texttt{KNOTTED}.  The respective numbers of crushing steps are $1,12,16$,
and all 29 labelled quotient comparisons agree.  Initial/shelled tetrahedron
counts are $20/5$, $60/48$, and $80/68$.  The two negative final sources have
10 and 13 tetrahedra, respectively, and one vertex each.

For these \emph{exact labelled final sources}, the audit establishes
\[
 h(\text{retained trefoil source})=3,
 \qquad
 h(\text{retained figure-eight source})=2.
\]
The trefoil lower bound uses positive root Q extreme rays with the pairwise
disjoint complete conflict sets
\[
 \{0,5\},\qquad \{1,2,3\},\qquad \{4,8\}.
\]
The figure-eight lower bound uses
\[
 \{0,1,2,3,12\},\qquad \{6,7,8,11\}.
\]
The complete negative certificates use only the split sets
$B=\{0,2,8\}$ and $B=\{8,12\}$, respectively, establishing the matching upper
bounds.  No completeness claim about enumerating positive root rays is
needed for either lower bound.

These lower witnesses originate in retained scratch discoveries, explicitly
separate from the fresh standard-surface enumeration.  All 19 retained root
rays are freshly replayed against the exact final source: nonnegative
integral original standard matching, zero canonical triangle minima,
positive formal Euler, primitive Q coordinates, and support nullity one in
the independently derived full-source Q equations.  The last condition proves
Q-ray extremality.  Incompatible coordinates here are formal vectors and are
never presented as embedded normal surfaces.  Source equality is literal
labelled equality; the checker refuses to transfer this evidence merely
because another triangulation has the same knot type or tetrahedron count.

Including unsuccessful candidates and restarted trees, the trefoil query
uses 64 oracle nodes, 124 LP calls, 1811 pivots, and 70 tested hitting subsets;
the figure-eight query uses 36 nodes, 114 LP calls, 5433 pivots, and 86 tested
subsets.  They consider seven and ten candidate resolving sets and retain
seven and eleven distinct learned conflicts.  Their anchor counts are 1644
and 733, with 1584 and 655 reuses, respectively.  Each prepares the fixed
source kernel once.  These operation counts are for the entire learning
process, not merely the final successful tree.  They demonstrate completion
within the common allowances of 256 nodes, 10000 pivots, 10000 anchors, and
10000 tested subsets; they establish neither a polynomial pivot bound nor a
bound on $h$ for a family of triangulations.

\subsection{Complete-query measurements and the practical tradeoff}

The affected regression suite passes all $160$ tests.  This includes complete
positive and negative diagram certificates, rejection of partial coverage,
strict limits and cancellation, and genuine search-driven crushing chains
that verify shared node, pivot, anchor and subset allowances across successors.
The shared oracle refactor also passes the existing source and restricted
DFS tests.  The default recognition dispatcher and default crushing sector
method remain unchanged.

The source-pinned paired experiment compares the current source-model DFS
with learning on the same three diagrams.  Both arms use $256$ nodes,
$10000$ pivots and $10000$ anchors; learning additionally permits $10000$
tested subsets.  Five randomized pairs and one shuffled DFS/DFS A/A pair
are recorded per case, with batches of one.  Construction, shelling, vertex
preprocessing, all failed candidates and subset filtering, LP discovery,
internal certificate replay, and full-result serialization are timed.
Imports, diagram objects and separate independent warmup replay are outside
the timer.  Every timed result has exactly its verified warmup's full output
hash.  Thus the learning measurements include unsuccessful candidate work.

\begin{center}
\small
\setlength{\tabcolsep}{3pt}
\begin{tabular}{lrrrrl}
\toprule
Case & DFS (ms) & Learning (ms) & DFS/learning & A/A & Outcomes \\
\midrule
circle & 8.661 & 8.661 & 1.008 & 0.998 & U/U \\
trefoil & 2887.522 & 2968.946 & 0.976 & 1.006 & K/K \\
figure\_eight & 11103.141 & 9302.810 & 1.194 & 1.155 & I/K \\
\bottomrule
\end{tabular}

\end{center}

Times are marginal median milliseconds; ratios are medians of paired ratios.
The outcome pairs are DFS/learning: U denotes \texttt{UNKNOT}, K denotes
\texttt{KNOTTED}, and I denotes \texttt{INCONCLUSIVE}.  Circle timings are
essentially unchanged.  On the trefoil the DFS and learning medians are
$2.888$ and $2.969$ seconds, with paired DFS/learning ratio $0.976$ and A/A
ratio $1.006$.  Learning gives a smaller final negative proof, $40$ nodes
instead of $64$, but is slightly slower in this sample.  Its $1811$ pivots
exceed DFS's $1670$, even though its anchor work is smaller.

On the figure-eight, DFS exhausts its anchor allowance without a completed
negative proof, while learning certifies knottedness with a $13$-node final
tree.  The measured $11.103$ and $9.303$ second medians compare these different
outcomes, so the paired ratio $1.194$ is not a speedup between two completed
recognitions.  Moreover, its A/A ratio is $1.155$, showing substantial timing
variation.  The stronger evidence here is finite completion under the shared
allowances, with the entire resulting certificate independently replayed.

The maintained package \path{fast/euler_transversal_research} retains the
fresh oracle audit, fixed lower-ray witnesses, paired runner and saved-evidence
checker.  The checker runs without Regina or discovery, validates all source
pins and warmup proofs, reconstructs the exact finite transversal bounds, and
recomputes the timing table.  The publication manifest also pins the article
and rebuilt PDF.  These results motivate an optional parameter-sensitive
search; they do not establish a uniform small resolving set or general
quasi-polynomial recognition.
